\chapter{BCR-ABL1 Genetic Testing}

\section*{What Is This Test?}
BCR-ABL1 genetic testing detects the presence and quantifies the level of the BCR-ABL1 fusion gene, the hallmark molecular abnormality of chronic myeloid leukemia (CML). The BCR-ABL1 fusion is formed by a reciprocal translocation between chromosomes 9 and 22 --- t(9;22)(q34;q11) --- producing the Philadelphia chromosome (Ph), a shortened chromosome 22. The fusion gene encodes a constitutively active tyrosine kinase (BCR-ABL1) that drives uncontrolled proliferation of myeloid cells. The breakpoint in the BCR gene determines the size of the fusion protein: the most common breakpoint (major breakpoint cluster region, M-bcr) produces the p210 BCR-ABL1 protein present in >95\% of CML; the minor breakpoint cluster region (m-bcr) produces the p190 protein, present in approximately 70\% of Philadelphia chromosome-positive acute lymphoblastic leukemia (Ph+ ALL) and in <5\% of CML; the mu-bcr produces the p230 protein, associated with neutrophilic-chronic myeloid leukemia. BCR-ABL1 testing is performed by quantitative reverse-transcription polymerase chain reaction (qRT-PCR) on peripheral blood or bone marrow, with results reported as the ratio of BCR-ABL1 transcripts to a reference gene (usually ABL1, BCR, or GUSB) on the International Scale (IS), expressed as \%IS. BCR-ABL1 monitoring is essential for assessing treatment response to tyrosine kinase inhibitors (TKIs --- imatinib, dasatinib, nilotinib, bosutinib, ponatinib, asciminib) and for detecting acquired resistance.

\section*{Alternative Names}
\begin{itemize}
  \item BCR-ABL1
  \item BCR-ABL
  \item Philadelphia Chromosome PCR
  \item BCR-ABL1 qRT-PCR
  \item BCR-ABL1 Transcript Detection
  \item BCR-ABL1 International Scale
  \item Minimal Residual Disease Monitoring in CML
  \item t(9
  \item 22) PCR
\end{itemize}
\section*{Principle}
BCR-ABL1 testing is performed by quantitative reverse-transcription PCR (qRT-PCR) on RNA extracted from peripheral blood leukocytes (or bone marrow aspirate). Whole blood is collected in PAXgene RNA tubes or EDTA tubes (for processing within 4 hours). RNA is extracted and reverse-transcribed to cDNA. Real-time quantitative PCR is performed using primers and TaqMan probes that specifically amplify the BCR-ABL1 fusion transcript (detecting both p210 [e13a2 and e14a2] and p190 [e1a2] isoforms) and a control reference gene (ABL1, BCR, or GUSB) for normalization. The BCR-ABL1/ABL1 (or BCR-ABL1/reference) ratio is calculated. Results are converted to the International Scale (IS) using a laboratory-specific conversion factor (CF) derived by comparison with the WHO International Genetic Reference Panel for BCR-ABL1. Results are reported as \%IS (BCR-ABL1 IS). A value of 10\% IS or 1\% IS represents 10\% or 1\% of the BCR-ABL1 transcript level found in standardized baseline samples. Key thresholds: complete cytogenetic response (CCyR) corresponds to approximately $\leq$1\% IS; major molecular response (MMR, MR3) is $\leq$0.1\% IS (a 3-log reduction from baseline); deep molecular response MR4 is $\leq$0.01\% IS; MR4.5 is $\leq$0.0032\% IS. The sensitivity of the assay (defined as the lower limit of detection) varies by laboratory but is typically approximately 0.001--0.01\% IS.

\section*{History}
The Philadelphia chromosome was discovered in 1960 by Peter Nowell (University of Pennsylvania) and David Hungerford (Fox Chase Cancer Center) in the bone marrow cells of patients with CML \cite{nowell1960minute} --- it was the first consistent chromosomal abnormality associated with a human cancer. In 1973, Janet Rowley used chromosome banding techniques to demonstrate that the Philadelphia chromosome results from a reciprocal translocation t(9;22) \cite{rowley1973consistent}. In 1984--1985, several groups cloned the BCR-ABL1 junction and demonstrated that the translocation fuses the ABL1 gene (chromosome 9) to the BCR gene (chromosome 22), producing a chimeric BCR-ABL1 fusion gene. The BCR-ABL1 protein was shown to have constitutive tyrosine kinase activity that is necessary and sufficient to cause CML in mouse models. In 1996, Brian Druker and colleagues demonstrated that the tyrosine kinase inhibitor imatinib (STI571, Gleevec) specifically inhibited the BCR-ABL1 kinase and selectively killed BCR-ABL1-positive cells \cite{druker1996effects}. Imatinib received accelerated FDA approval in 2001 after the landmark IRIS trial demonstrated unprecedented efficacy. The International Scale for BCR-ABL1 standardization was established in 2003--2007 to harmonize results across laboratories and clinical trials \cite{branford2008desirable}. Second-generation (dasatinib, nilotinib, bosutinib) and third-generation (ponatinib) TKIs were introduced to overcome resistance. Asciminib, the first allosteric BCR-ABL1 inhibitor (STAMP inhibitor, Specifically Targeting the ABL Myristoyl Pocket), was approved in 2021 for TKI-resistant disease.

\section*{Why This Test Is Ordered}
Diagnosis of CML and Ph+ ALL: initial BCR-ABL1 qRT-PCR is performed on peripheral blood or bone marrow at diagnosis to confirm the presence of the fusion transcript and to establish the baseline BCR-ABL1 level. Molecular monitoring of CML during TKI therapy: qRT-PCR is performed every 3 months after initiating TKI therapy. The 2020 European LeukemiaNet (ELN) and NCCN guidelines define treatment milestones \cite{hochhaus2020european}: BCR-ABL1 $\leq$10\% IS at 3 months (optimal), $\leq$1\% IS (CCyR equivalent) at 6 months (optimal), $\leq$0.1\% IS (MMR/MR3) at 12 months (optimal). Failure to achieve these milestones constitutes TKI failure and warrants a change in therapy. Detection of TKI resistance: a confirmed rise in BCR-ABL1 levels (defined as loss of MMR on two consecutive tests) may indicate acquired resistance from BCR-ABL1 kinase domain mutations. Kinase domain mutational analysis (Sanger sequencing or NGS of the ABL1 kinase domain) should be performed. Treatment-free remission (TFR): patients who have achieved sustained deep molecular response (MR4 or MR4.5 for $\geq$2 years) may be candidates for TKI discontinuation with serial molecular monitoring. Detection of minimal residual disease in Ph+ ALL.

\section*{Primary vs Secondary Test}
This is a primary diagnostic test for CML and Ph+ ALL and a primary monitoring test for CML treatment response. BCR-ABL1 qRT-PCR is essential for the diagnosis, risk stratification, and management of CML, replacing bone marrow cytogenetics for routine monitoring after CCyR is achieved.

\section*{How This Test Is Performed}
A blood sample (5--10 mL) is collected in a PAXgene RNA tube or an EDTA (lavender-top) tube. For bone marrow testing, 1--2 mL of bone marrow aspirate is collected. RNA is extracted and analyzed by qRT-PCR. Peripheral blood is the preferred specimen for routine CML monitoring because BCR-ABL1 levels in peripheral blood closely correlate with bone marrow levels.

\section*{What Preparation Is Needed}
No special preparation is required. Fasting is not necessary. The blood specimen should be processed within 4 hours if collected in an EDTA tube (to prevent RNA degradation); PAXgene tubes stabilize RNA for up to 3 days at room temperature. Serial monitoring should ideally be performed in the same laboratory using the same assay to minimize variability; if a change in laboratory is necessary, at least one concurrent sample should be tested in both laboratories to establish correlation.

\section*{Contraindications and Precautions}
\begin{itemize}
  \item No absolute contraindications. The most important consideration is specimen handling: RNA is labile and degrades rapidly. Delayed processing, repeated freeze-thaw cycles, or improper storage of RNA degrades BCR-ABL1 transcripts and can produce false-negative results. The International Scale (IS) conversion factor must have been validated for the specific assay and laboratory. Results reported as \%IS can only be generated by laboratories that have established and validated their conversion factor. BCR-ABL1 p190 (e1a2) in CML is rare (<5\%) and may not be detected by assays designed for p210 (e13a2/e14a2) alone; laboratories should verify that their assay detects all common BCR-ABL1 transcripts. Atypical BCR-ABL1 transcripts (e19a2, e6a2, e8a2) exist and may require specialized testing. A rise in BCR-ABL1 transcript levels of $\geq$5-fold (confirmed on a repeat sample) with loss of MMR should prompt testing for BCR-ABL1 kinase domain mutations, as this often indicates emerging TKI resistance.
\end{itemize}
\section*{Risks}
Minimal risk from venipuncture. Slight pain or bruising at the site. Rare: hematoma.

\section*{What Happens After the Test}
Results are typically available within 3--7 days. BCR-ABL1 levels are compared with previous results and validated against ELN/NCCN milestones. A BCR-ABL1 level that fails to decrease appropriately (or that rises after prior decline) prompts evaluation for: medication non-adherence (the most common cause of apparent TKI failure), drug-drug interactions (imatinib is metabolized by CYP3A4; inducers like rifampin and St. John's Wort reduce TKI levels), BCR-ABL1 kinase domain mutations (testing for ABL1 mutations), and clonal evolution (cytogenetic analysis for additional chromosomal abnormalities). Loss of MMR on two consecutive samples with confirmed TKI resistance warrants consideration of a change in TKI therapy.

\section*{Understanding Results}

\subsection*{BCR-ABL1 Not Detected (Undetectable Molecular Residual Disease)}
BCR-ABL1 transcript not detected by the assay (the assay's sensitivity limit, typically MR4.5 [0.0032\% IS] or MR5 [0.001\% IS]). A patient with undetectable BCR-ABL1 for $\geq$2 years (confirmed MR4.5 for $\geq$2 years) on TKI therapy is potentially eligible for a trial of treatment-free remission (TFR). Approximately 40--50\% of patients who discontinue TKIs in sustained deep molecular response maintain treatment-free remission; 50--60\% experience molecular relapse (detectable BCR-ABL1) and must restart TKIs. Relapse after TKI discontinuation is typically sensitive to re-initiation of the same TKI.

\subsection*{BCR-ABL1 Detected (Quantifiable Transcripts)}
At diagnosis: BCR-ABL1 IS is typically 30--100\% IS. At 3 months of TKI therapy: BCR-ABL1 $\leq$10\% IS is the optimal response; >10\% IS at 3 months is a warning response that requires more frequent monitoring (molecular testing every 1--3 months rather than every 3 months) but does not automatically mandate a change in therapy. At 6 months: BCR-ABL1 $\leq$1\% IS (CCyR equivalent) is optimal; >1--10\% IS is a warning response; >10\% IS is failure and requires a change in TKI. At 12 months: BCR-ABL1 $\leq$0.1\% IS (MMR/MR3) is optimal; >0.1--1\% IS is a warning response (failure per 2020 ELN if >1\% IS at 12 months). Loss of MMR (confirmed BCR-ABL1 >0.1\% IS on two consecutive samples after previously achieving MMR): this is a TKI failure and requires BCR-ABL1 kinase domain mutation testing. A confirmed rise of $\geq$5-fold in BCR-ABL1 from the prior measurement should also prompt kinase domain mutation analysis. In patients with the T315I kinase domain mutation, third-generation TKI therapy is indicated: ponatinib is active against T315I; asciminib is also active against T315I per updated data.

\section*{Normal Results}
BCR-ABL1 NOT DETECTED (undetectable). In healthy individuals (and individuals without CML or Ph+ ALL), no BCR-ABL1 transcript is present; the result is ``BCR-ABL1 not detected.'' Low-level BCR-ABL1 transcripts are detectable by highly sensitive PCR in a small proportion of healthy individuals (approximately 1 in 10,000 cells) due to the occurrence of random breakage events at the BCR and ABL1 loci; these levels are extremely low and do not meet criteria for CML. The following milestones define treatment response in CML: 3 months: $\leq$10\% IS (optimal), >10\% IS (warning). 6 months: $\leq$1\% IS (optimal), >1--10\% IS (warning), >10\% IS (failure). 12 months: $\leq$0.1\% IS/MMR (optimal), >0.1--1\% IS (warning per ELN 2013; failure per ELN 2020), >1\% IS (failure). Any time: loss of MMR on consecutive tests (failure). Deep molecular response: MR4 ($\leq$0.01\% IS), MR4.5 ($\leq$0.0032\% IS), MR5 ($\leq$0.001\% IS).

\section*{Follow-up Tests}
BCR-ABL1 kinase domain mutation analysis (Sanger sequencing or NGS) when TKI failure or loss of MMR occurs. Bone marrow cytogenetics (karyotype) at diagnosis and if there is concern for clonal evolution, myelodysplasia, or blast crisis transformation. Complete blood count every 1--2 weeks during the first 3 months of TKI therapy, and then every 1--3 months. Bone marrow biopsy and aspiration if the patient does not achieve hematologic or cytogenetic milestones. Pharmacokinetic testing (imatinib trough levels) if non-adherence or drug interactions are suspected. ABL1 mutation-specific therapy: ponatinib for T315I; alternative TKI for other kinase domain mutations. Switch to asciminib if multiple TKIs have failed or if intolerance prevents continued use.

\section*{References}
\bibliographystyle{unsrtnat}
\bibliography{references}

\clearpage
\section*{Regulatory Status and Authorization Date}
This chapter describes a test category, not a specific commercial product. FDA, CDSCO, EU IVDR, and MHRA approval or registration dates therefore cannot be assigned without the assay manufacturer, product name, instrument, and intended use. For laboratory-developed tests, an individual agency approval date may not exist. See the Preface for the interpretation of regulatory status.

\clearpage
